\section{Handlungsoptionen des Anlegers}
\label{sec:optionen}

Der Anleger, nicht das Unternehmen, entscheidet hier. Zur Wahl stehen drei Optionen, die
sich gegenseitig ausschlie{\ss}en:

\begin{description}
\item[Option A] Kauf zum heutigen Kurs von \je{\OptionAKurs}{EUR}.
\item[Option B] Kauf erst zum Schwellenkurs von \je{\OptionBKurs}{EUR}, mit dem Risiko,
  dass dieser Kurs nicht erreicht wird.
\item[Option C] Verzicht auf GEA, Anlage desselben Betrags in den Weltindex zur H\"urde von
  \Pz{\Huerde}.
\end{description}

\optionentabelle

\subsection{Option A: Kauf heute}

\begin{table}[htbp]\centering\small
\caption{Bewertungsraster Option A.}\label{tab:raster-a}
\begin{tabularx}{\linewidth}{@{}lX@{}}\toprule
Faktor & Einsch\"atzung\\\midrule
Investitionsbedarf & Der volle Kaufpreis von \je{\OptionAKurs}{EUR} je Aktie ist sofort
f\"allig; Nachschuss entf\"allt.\\
Kapitalbindung & Vollst\"andig und unbefristet. Der Buchwert je Aktie betr\"agt
\je{\BuchwertJeAktie}{EUR}, also \KBV{} mal weniger als der Kaufpreis; im Liquidationsfall
w\"are nur dieser Teil gedeckt.\\
Liquidit\"atsbedarf (j\"ahrlich) & Keiner. Die Anlage liefert einen Zufluss:
\je{\Dividende}{EUR} je Aktie f\"ur 2025.\\
Zeithorizont bis Return & Bei konstantem Zufluss und Buchwert als Endwert erreicht die
Anlage die H\"urde auf keinem der drei Horizonte (Tabelle~\ref{tab:optionen}).\\
Risiko & Das Gesch\"aft ist stetig und ohne Nettoschulden; das Risiko liegt nicht im
Unternehmen, sondern im Preis.\\
Potenzial & Voll am Wachstum beteiligt. Erreicht GEA die Ziele von Mission 30, w\"achst das
EBITDA um \MissionImplizit{} p.\,a.\\
\bottomrule\end{tabularx}\end{table}

\bk{\emph{Dauer der Umsetzung:} sofort. \emph{Rentabilit\"at:} zum heutigen Kurs negativ
gegen\"uber der H\"urde, \OptionAIrrMed{} auf zehn Jahre im Basisfall.
\emph{Liquidit\"at:} unproblematisch, die Position bindet kein weiteres Geld und
aussch\"uttet. \emph{Wirkung auf das Kerngesch\"aft des Anlegers:} Wer GEA kauft, kauft
einen stetigen Zahler zu einem Preis, der Wachstum voraussetzt.}

\subsection{Option B: Kauf zum Schwellenkurs}

\begin{table}[htbp]\centering\small
\caption{Bewertungsraster Option B.}\label{tab:raster-b}
\begin{tabularx}{\linewidth}{@{}lX@{}}\toprule
Faktor & Einsch\"atzung\\\midrule
Investitionsbedarf & \je{\OptionBKurs}{EUR} je Aktie, gut ein Viertel des heutigen Kurses.\\
Kapitalbindung & Bis zum Erreichen des Kurses null, danach wie Option~A.\\
Liquidit\"atsbedarf (j\"ahrlich) & Keiner; der Betrag bleibt bis dahin anlegbar.\\
Zeithorizont bis Return & Unbestimmt. Der Kurs lag zuletzt 2018 in dieser
Gr\"o{\ss}enordnung; der Zeitpunkt ist nicht planbar.\\
Risiko & Das Risiko der Option ist nicht der Verlust, sondern das Ausbleiben:
\je{\OptionBKurs}{EUR} setzen einen Kursr\"uckgang voraus, den kein belegter Umstand
erwarten l\"asst.\\
Potenzial & Bei Zustandekommen \OptionBIrrMed{} auf zehn Jahre, also \"uber der H\"urde.\\
\bottomrule\end{tabularx}\end{table}

\bk{\emph{Dauer der Umsetzung:} unbestimmt. \emph{Rentabilit\"at:} die h\"ochste der drei
Optionen, aber nur unter einer Bedingung, die der Anleger nicht beeinflusst.
\emph{Liquidit\"at:} das Geld bleibt bis zum Kauf verf\"ugbar und kann zwischenzeitlich
die H\"urde verdienen. \emph{Wirkung:} Diese Option ist in der Praxis Option~C mit einer
Kauforder.}

\subsection{Option C: Verzicht und Indexanlage}

\begin{table}[htbp]\centering\small
\caption{Bewertungsraster Option C.}\label{tab:raster-c}
\begin{tabularx}{\linewidth}{@{}lX@{}}\toprule
Faktor & Einsch\"atzung\\\midrule
Investitionsbedarf & Derselbe Betrag, anderer Gegenstand.\\
Kapitalbindung & Vollst\"andig, aber \"uber hunderte Titel gestreut.\\
Liquidit\"atsbedarf (j\"ahrlich) & Keiner.\\
Zeithorizont bis Return & Die H\"urde ist die Rendite dieser Option; sie wird
definitionsgem\"a{\ss} auf jedem Horizont erreicht.\\
Risiko & Marktrisiko ohne Einzeltitelrisiko. Die \Pz{\Huerde} sind eine Vergangenheit,
keine Zusage.\\
Potenzial & \OptionCIrr{} p.\,a.\ Kein Titel dieses Berichts \"ubertrifft das zum heutigen
Kurs.\\
\bottomrule\end{tabularx}\end{table}

\bk{Option~C ist keine Verlegenheitsl\"osung, sondern die Messlatte. Eine Einzelanlage muss
sie schlagen, sonst tr\"agt der Anleger Einzeltitelrisiko ohne Gegenleistung. Der Vergleich
ist allerdings nicht ganz fair zugunsten von C: Die \Pz{\Huerde} sind eine
Zehnjahresrendite, die stark von einer einzigen Periode getragen ist; \Pz{\HuerdeLang}
seit Auflage ist der vorsichtigere Ma{\ss}stab, und er hebt den Schwellenkurs von
Option~B auf \je{\SchwelleMEDZehnLang}{EUR}.}
